% Chapter 3, typeset from lectures/L03: the README is the opener and the closing review, and each
% part of Appendix A is a section, in the same order and with the same subsections. Appendix B,
% the exercises, is the last section.

\chapter{Sequential Logic}
\label{c3:ch}

{\large\itshape Storing state: registers, D flip-flops, and clocking.\par}

\begin{chapterbox}{In this chapter}
\begin{itemize}
  \item Why combinational logic alone cannot remember anything, and what ``state'' means in a
    circuit.
  \item The D latch and the D flip-flop, and why synchronous designs are built from the flip-flop.
  \item Clocking: period, frequency, and rising versus falling edges.
  \item The synchronous process template in VHDL, and when a signal assignment actually takes
    effect.
  \item Edge detection: turning a level into a one-cycle pulse.
  \item \code{led_toggle} on a DE0-CV: a button press that toggles an LED and holds it.
  \item Registers as flip-flops in parallel, in \secref{c3:a5}, and \code{d_flip_flop} on its own,
    in \secref{c3:a8}.
\end{itemize}
\end{chapterbox}

The course this book comes from builds this chapter live, one circuit growing out of the last. It
starts with a D latch in CircuitVerse, cross-coupled gates that are transparent while enabled, and
makes a D flip-flop from a pair of them: the same circuit, now changing only on an edge. Then comes
an edge-detected LED toggle, also in CircuitVerse, with one button: one flip-flop remembers the
previous button state, a second holds the LED, and a falling edge on the button toggles it. The
same circuit is then written in VHDL for two buttons at once, the synchronous process template
twice over, with the single wires of the drawing widened into 2-bit vectors so that one process
serves both buttons. That widening is the whole difference between the drawing and the module.
Finally \code{led_toggle} goes through Quartus onto the DE0-CV: press the button, and watch the LED
hold its new state until the next press. \Secref{c3:a9} writes up that last part.

Two predictions are worth making before each simulation runs:
\begin{itemize}
  \item for the latch, what \code{Q} does when \code{enable} goes low while \code{D} is still
    changing.
  \item for the flip-flop, what \code{Q} does when \code{D} changes \emph{between} two clock edges.
\end{itemize}
Exercises~\ref{c3:ex:1} and~\ref{c3:ex:2} are the latch and the flip-flop, which is where those two
predictions are worth confirming rather than taking on trust. The module \code{d_flip_flop} is also
written up on its own, in \secref{c3:a8}, with a reference testbench. It is for reading rather than
presenting: the course builds the flip-flop as part of the toggle circuit instead, which is the same
module arriving with something to do.

Read the sections that follow before you start the exercises. This is the first chapter whose
exercises need a testbench for a clocked design, so \secref{c2:app:c} is worth rereading before you
start \code{register4} in Exercise~\ref{c3:ex:3}. The board work is shown in
Appendix~\ref{app:quartus} for anyone who has a DE0-CV; no exercise needs one.

% Sections 3.1 to 3.9, from lectures/L03/appendix/a_flip_flops_and_registers.md: one section per
% part, A.1 to A.9.
%
% The four block drawings are drawn with box-drawing characters, so they are textdrawing blocks.

\section{From Combinational to Sequential Logic}
\label{c3:a1}
Every circuit in Chapters~\ref{c1:ch} and~\ref{c2:ch} was \textbf{combinational}: its output is a
pure function of its present inputs, so the same inputs always give the same output, and there is
no notion of ``before'' or ``after'' inside the circuit.

Most useful systems need more. A counter must remember its count. A register must hold a value after
the signal that produced it has changed. An LED toggled by a button must remember whether it is on,
since that is not encoded in the button's momentary state at all.

All of these need \textbf{memory}: an output that depends on the circuit's own \emph{previous}
output. Sequential logic is combinational logic with part of the output fed back into the input, and
that \textbf{feedback loop} is what turns a stateless gate network into a circuit with state. This
chapter builds the primitive that makes controlled feedback practical, the D flip-flop, then the two
things you build directly out of it: registers and edge detectors.

\section{The D Latch}
\label{c3:a2}
The simplest storage element is the \textbf{D latch}: two cross-coupled gates, gated by an
\code{enable} input, storing a single bit.

\needspace{9\baselineskip}
\begin{textdrawing}
        ┌────────────┐
   D ──►│            │──► Q
        │  D LATCH   │
enable─►│            │──► Qn
        └────────────┘
\end{textdrawing}

\noindent\code{D} is the value to store, \code{enable} controls whether the latch is \emph{open}
(transparent) or \emph{closed} (holding), \code{Q} is the stored value and \code{Qn} is always its
inverse.

Each output feeds back into the other's equation, which is exactly \secref{c3:a1}'s feedback loop:
\[ Q = (D' \cdot \bool{enable} + \bool{Qn})' \]
\[ \bool{Qn} = (D \cdot \bool{enable} + Q)' \]

\begin{booktable}{@{}p{3.6em}p{13em}L@{}}
  \thead{\code{enable}} & \thead{Behaviour} & \thead{Result} \\
  \midrule
  \code{1} & latch is \textbf{open} (transparent) & \code{Q = D}, \code{Qn = D'}; output follows
    the input immediately \\
  \code{0} & latch is \textbf{closed} (locked) & \code{Q} and \code{Qn} hold whatever they last
    had \\
\end{booktable}

\noindent\textbf{Build it in CircuitVerse} from one NOT gate, two AND gates and two NOR gates (a NOR
is an OR followed by a NOT). Cross-couple the feedback so the NOR producing \code{Q} feeds the NOR on
the \code{Qn} side and vice versa. Drive \code{D} and \code{enable} from input switches and watch
\code{Q} and \code{Qn} on LEDs.

Then test it:
\begin{itemize}
  \item With \code{enable = 1}, change \code{D} and confirm \code{Q} follows immediately. The latch
    is \textbf{transparent}.
  \item With \code{enable = 0}, confirm \code{Q} is unchanged whatever \code{D} does. The latch is
    \textbf{locked}.
  \item Cycle \code{D} a few times while toggling \code{enable}, and confirm the latch always freezes
    whatever \code{Q} held at the instant \code{enable} fell. That is what
    \textbf{level-sensitive} means: transparent for the entire time \code{enable} is high, not just
    at one instant.
\end{itemize}
That last point is also the latch's fundamental weakness. Because it is transparent for the whole
time \code{enable = 1}, any glitch on \code{D} during that window passes straight through to
\code{Q}. In a large synchronous system, guaranteeing stable data throughout an entire clock phase
is far harder than guaranteeing it around a single instant. That is the problem the D flip-flop
solves.

\section{The D Flip-Flop}
\label{c3:a3}
A \textbf{D flip-flop} stores a bit the same way a latch does, but is \textbf{edge-triggered} rather
than level-sensitive: it samples \code{D} only at the \emph{instant} the clock transitions, not for
the whole time the clock is high or low.

Feed the same \code{D} and the same clock to both and the difference is immediate
(Figure~\ref{c3:fig:latch_vs_flipflop}).

\bookfigure{latch_vs_flipflop}{A latch and a flip-flop fed the same \code{D} and clock: the latch
follows \code{D} throughout every clock-high window, while the flip-flop changes only on the rising
edges.}{c3:fig:latch_vs_flipflop}

\noindent The latch changes three times: twice because \code{D} moved while the clock was high, and
once because the clock went high onto a \code{D} that had already changed. The flip-flop changes
once, because only one of the four rising edges found \code{D} different from the value \code{Q} was
already holding. Both are behaving correctly; they are answering different questions.

\needspace{9\baselineskip}
\begin{textdrawing}
        ┌───────────────┐
   D ──►│               │──► Q
clock ─►│  D FLIP-FLOP  │
enable─►│ (rising-edge) │──► Qn
reset_n►│               │
        └───────────────┘
\end{textdrawing}

\noindent\code{clock} is the timing reference (\secref{c3:a4}). \code{enable} gates whether the
flip-flop accepts a new value \emph{at the triggering edge}. \code{reset_n} is an asynchronous,
active-low reset: whenever it is \code{0}, \code{Q} is forced to \code{0} immediately, regardless of
the clock.

\begin{booktable}{@{}>{\raggedright\arraybackslash}p{17em}L@{}}
  \thead{Condition} & \thead{Result} \\
  \midrule
  \code{reset_n = 0} & \code{Q = 0}, \code{Qn = 1}, immediately (asynchronous, ignores the clock
    entirely) \\
  \code{reset_n = 1}, rising edge on \code{clock}, \code{enable = 1} & \code{Q = D},
    \code{Qn = D'}, sampled at that instant \\
  \code{reset_n = 1}, rising edge on \code{clock}, \code{enable = 0} & \code{Q} holds its previous
    value \\
  \code{reset_n = 1}, no rising edge on \code{clock} & \code{Q} holds its previous value, no matter
    what \code{D} does \\
\end{booktable}

\noindent That last row is worth seeing rather than reading (Figure~\ref{c3:fig:dff_timing}).

\bookfigure{dff_timing}{A flip-flop between edges: \code{D} rises mid-cycle and is captured at the
next rising edge, then dips and recovers entirely between two edges without \code{Q} ever
moving.}{c3:fig:dff_timing}

\noindent\code{D} dips low and comes back inside a single clock period, and \code{Q} never moves,
because no edge happened while it was low. A latch in the same position would have followed it down
and back up.

Internally, an edge-triggered flip-flop is two D latches in series, a \textbf{master-slave} pair. One
is transparent while the clock is low and captures \code{D}; the other is transparent while the clock
is high and passes the captured value to \code{Q}. The net effect is that \code{Q} changes only at
the instant the clock goes low to high, never while it sits at either level. That external behaviour
is what VHDL's \code{rising_edge()} models (\secref{c3:a6}).

\textbf{Build it in CircuitVerse} by extending your \secref{c3:a2} latch:
\begin{itemize}
  \item Duplicate the latch so two sit in series: the \emph{master} feeds the \emph{slave}, whose
    \code{Q} is the flip-flop's output.
  \item Add \code{clock}, \code{reset_n}, \code{D} and \code{enable} as inputs.
  \item Drive the two latches' transparency from the clock and its inverse. To keep the two meanings
    apart, call that internal per-latch control \code{gate}: the master's \code{gate} is
    \code{clock'} and the slave's is \code{clock}. The flip-flop's external \code{enable} is a
    different signal and is wired separately, below.
  \item Wire the external \code{enable} as a 2-to-1 mux (\secref{c2:a4}) on the master's data input:
    when \code{enable = 1} the mux passes \code{D}, and when \code{enable = 0} it passes the
    flip-flop's own \code{Q} back in, so the edge re-captures the value already stored. Gate the
    \emph{data}, never the clock.
  \item Add the reset to \textbf{both} latches. Forcing \code{Q} low takes two changes in each
    latch, not one:
    \begin{itemize}
      \item feed an active-high \code{reset} (\code{reset_n} through a NOT) into the NOR producing
        \code{Q}, forcing \code{Q} to \code{0}.
      \item AND \code{reset_n} into that same latch's \code{D AND gate} term, which frees \code{Qn}
        to go high.
    \end{itemize}
  \item Set the clock period to something slow enough to watch, e.g.\ \code{1000 ms}.
\end{itemize}
With the reset wired in, each latch's equations become:
\[ Q = (D' \cdot \bool{gate} + \bool{Qn} + \bool{reset})' \]
\[ \bool{Qn} = (D \cdot \bool{gate} \cdot \bool{reset\_n} + Q)' \]
A tempting shortcut is to leave the \code{Qn} equation alone and break the loop instead, ANDing
\code{reset_n} onto the \code{Qn} wire feeding \code{Q}'s NOR. That is not enough. With
\code{gate = 1} and \code{D = 1} it leaves \code{Qn} at \code{0}, which is what holds \code{Q} at
\code{1}, so the reset silently does nothing in exactly the case you are most likely to test it in.
Both changes, in both latches, or neither.

Confirm that \code{Q} updates only at the instant the clock rises, holding steady between edges even
if \code{D} changes, and that pulling \code{reset_n} low forces \code{Q} to \code{0} immediately,
with releasing it changing nothing until the next rising edge.

\section{Clocking}
\label{c3:a4}
The \textbf{clock} is a periodic square wave, read by every flip-flop in a synchronous circuit from
the same source. Two numbers describe it: the \textbf{period} \code{T}, the time for one full cycle,
and the \textbf{frequency} \code{f = 1/T} in Hz.

Each period has exactly two \textbf{edges}: a \textbf{rising edge} from \code{0} to \code{1}, and a
\textbf{falling edge} from \code{1} to \code{0}. A synchronous design picks one of them, almost
always the rising edge, and updates every flip-flop on that edge and only that edge. Using a single
edge everywhere is what makes a large circuit's timing analyzable: every flip-flop changes at the
same well-defined instants, so you can reason one clock cycle at a time rather than tracking every
signal continuously.

Simulated and real clocks run on wildly different timescales, and the logic is identical either way:
\begin{itemize}
  \item In CircuitVerse, set a period you can watch, e.g.\ \code{1000 ms}.
  \item On the DE0-CV the system clock is a fixed \code{50 MHz}, a period of \code{20 ns}. Nobody
    perceives individual edges at that rate, so you reason in terms of ``on the next rising edge'',
    never in wall-clock time.
\end{itemize}

\section{Registers: Flip-Flops in Parallel}
\label{c3:a5}
A \textbf{register} stores more than one bit: N D flip-flops side by side, one per bit, all sharing
the same \code{clock}, \code{reset_n} and \code{enable}.

\needspace{9\baselineskip}
\begin{textdrawing}
D(0) ──►[D FF]──► Q(0)
D(1) ──►[D FF]──► Q(1)
D(2) ──►[D FF]──► Q(2)
   ⋮        ⋮
D(N-1)──►[D FF]──► Q(N-1)

  (clock, reset_n, enable shared by every flip-flop above)
\end{textdrawing}

\noindent That is the entire idea: an 8-bit register is eight D flip-flops updating on the same clock
edge.

In VHDL you rarely instantiate N flip-flops explicitly. Declare one \code{std_logic_vector} and
assign the whole thing inside a single synchronous process (\secref{c3:a6}): every bit synthesizes
to its own flip-flop, all sharing the process's clock and reset.

\needspace{14\baselineskip}
\begin{vhdlcode}
signal q: std_logic_vector(3 downto 0);
...
process (clock, reset_n) is
begin
    if (reset_n = '0') then
        q <= "0000";
    elsif (rising_edge(clock)) then
        if (enable = '1') then
            q <= d; -- d is a std_logic_vector(3 downto 0) input
        end if;
    end if;
end process;
\end{vhdlcode}

\noindent Four bits, one process, four flip-flops in hardware: no different in principle from four
\code{d_flip_flop} instances wired in parallel.

\section{The Synchronous Process Template in VHDL}
\label{c3:a6}
Every synchronous element in this chapter, the flip-flop, the register and the edge detector, uses
the same shape. Without a reset, a single D flip-flop is:

\begin{vhdlcode}
process (clock) is
begin
    if (rising_edge(clock)) then
        q <= d;
    end if;
end process;
\end{vhdlcode}

\noindent\code{q} takes \code{d}'s value at every rising edge and holds it the rest of the time. That
one line inside the \code{if} is enough for synthesis to infer a real D flip-flop.

\subsection{Rules}
\label{c3:rules}
\begin{itemize}
  \item \textbf{\code{clock} is the only signal in the sensitivity list}, unless there is an
    asynchronous reset. The process should run only when the clock changes; it has no reason to
    react to anything else.
  \item A process reacts to \emph{any} change on a sensitivity-list signal, not just rising ones, so
    \code{rising_edge()} filters for the rising edge specifically. All logic that updates
    synchronously goes inside that \code{if}.
  \item Every signal assigned inside such a process synthesizes to a flip-flop, one bit each.
\end{itemize}

\subsection{With an asynchronous reset}
\label{c3:asyncreset}
An asynchronous reset must also appear in the sensitivity list, since it can be asserted
independently of the clock, and it takes priority when active.

\needspace{10\baselineskip}
\begin{vhdlcode}
process (clock, reset_n) is
begin
    if (reset_n = '0') then
        q <= '0';
    elsif (rising_edge(clock)) then
        q <= d;
    end if;
end process;
\end{vhdlcode}

\noindent Read it in priority order: check the reset first, and only if it is inactive check for a
rising edge. If neither holds (a falling edge, or \code{reset_n} going from \code{0} back to
\code{1}, which also re-triggers the process), nothing is assigned and \code{q} keeps its previous
value.

\subsection{When the assignment actually happens}
\label{c3:scheduling}
One rule sits underneath all of the above, and it is the single place where \code{<=} behaves unlike
assignment in any language you have written software in.

A signal assignment does not take effect when its line runs. It \textbf{schedules} a value, applied
only once the process finishes this pass and suspends. Until then, every read of that signal returns
the value it held \emph{before the pass started}, however many assignments have already executed.

That sounds like a technicality. It is the reason the template above is a flip-flop rather than a
wire:

\begin{vhdlcode}
process (clock) is
begin
    if (rising_edge(clock)) then
        b <= a;
        c <= b;
    end if;
end process;
\end{vhdlcode}

\noindent Read as software, \code{b} and \code{c} both end up holding \code{a}. In hardware they do
not. On each edge \code{b} is scheduled to take \code{a}'s value, and \code{c} is scheduled to take
the value \code{b} held \emph{coming into this edge}, not the one just scheduled onto it. The result
is two flip-flops in a chain: \code{a} reaches \code{b} after one edge and \code{c} after two. Swap
the lines and nothing changes, because neither assignment can observe the other's result.

Two consequences worth carrying forward:
\begin{itemize}
  \item \textbf{Order does not matter} between assignments to \emph{different} signals in the same
    clocked process. They all read the old values and all take effect together. That is the opposite
    of the sequential reading the word ``process'' invites, and it is why a chain of flip-flops can
    be written in any order.
  \item \textbf{Assign the same signal twice in one pass and only the last assignment survives.} The
    earlier ones are discarded without ever reaching the signal.
\end{itemize}
\Secref{c3:a7}'s edge detector is exactly the two-element chain above, and works for exactly this
reason. Chapter~\ref{c5:ch} returns to this rule from the other side, with a \code{variable}, which
updates immediately and so behaves the way the software reading expects.

\section{Edge Detection}
\label{c3:a7}
A flip-flop's other everyday job is detecting the \emph{instant} a signal changes rather than reading
its level: store the signal's previous value in a flip-flop, then compare it against the current
value with a gate.
\begin{itemize}
  \item \textbf{Rising-edge detection} (\code{0} to \code{1}):
    \[ \bool{edge} = \bool{current} \cdot \bool{previous}' \]
  \item \textbf{Falling-edge detection} (\code{1} to \code{0}):
    \[ \bool{edge} = \bool{current}' \cdot \bool{previous} \]
\end{itemize}
The rising-edge form as a circuit; swap which side is inverted for a falling edge, which is the form
\secref{c3:a9}'s active-low button needs:

\needspace{9\baselineskip}
\begin{textdrawing}
signal ──┬───────────────────────────► AND ──► edge
         │                              ▲
         └──►[D FF]── previous ──(NOT)──┘
                ▲
           clock, reset_n
\end{textdrawing}

\noindent\code{edge} is high for exactly one clock cycle, the one immediately after the transition,
because \code{previous} only catches up on the \emph{next} edge. This is the mechanism behind ``do
something once per button press'' rather than ``do something every clock cycle the button is held'',
which is exactly what \secref{c3:a9} solves.

\section{Worked Example: A D Flip-Flop in VHDL}
\label{c3:a8}
\file{lectures/L03/d_flip_flop/d_flip_flop.vhd} implements \secref{c3:a3} and \secref{c3:a6}
combined: an enable gating whether a rising edge captures \code{d}, and an asynchronous active-low
reset.

\begin{vhdlcode}
entity d_flip_flop is
    port(clock, reset_n, d, enable: in std_logic;
         q, q_n                   : out std_logic);
end entity;
\end{vhdlcode}

\noindent Its single process is the reset-plus-enable variant of \secref{c3:a6}'s template: reset
takes priority, then a rising edge captures \code{d} into an internal signal \code{q_s} if
\code{enable = '1'}. \code{q_s} exists because an \emph{output} port cannot be read back inside the
same architecture; \code{q} and \code{q_n} are both driven combinationally from it, outside the
process. This is the whole of \file{d_flip_flop.vhd}:

\vhdlfile{lectures/L03/d_flip_flop/d_flip_flop.vhd}

To try it on hardware, assign \code{clock} to the \code{50 MHz} system clock, \code{reset_n} to a
pushbutton, \code{d} and \code{enable} to two slide switches, and \code{q}/\code{q_n} to two LEDs.
Then flip \code{enable} on and change \code{d}, and the LED follows; flip \code{enable} off and change
\code{d}, and it does not.

There is no clock edge you can perceive at \code{50 MHz}, so what you are really confirming is that
the logic still holds once you can no longer see the individual edges. The flip-flop looks like it
responds to your switch; in fact it responded to one of fifty million clock edges that second, and
the switch only decided which value that edge captured.

\section{Worked Example: Edge-Detected LED Toggle}
\label{c3:a9}
\file{lectures/L03/led_toggle/led_toggle.vhd} combines \secref{c3:a5}, \secref{c3:a6} and
\secref{c3:a7} into one small system: two independent push buttons, each toggling its own LED
exactly once per press.

\begin{vhdlcode}
entity led_toggle is
    port(clock, reset_n: in std_logic;
         button_n      : in std_logic_vector(1 downto 0);
         led           : out std_logic_vector(1 downto 0));
end entity;
\end{vhdlcode}

\noindent Three 2-bit signals, two held in clocked processes and one computed between them:
\begin{itemize}
  \item \code{button_prev_n} is a 2-bit register (\secref{c3:a5}) holding last cycle's
    \code{button_n}, reset to \code{"11"} since the buttons are active-low.
  \item \code{button_edge <= (not button_n) and button_prev_n} is \secref{c3:a7}'s falling-edge
    detector applied to both bits at once, high for one cycle exactly when a button goes from
    released to pressed.
  \item \code{led_s} is a second 2-bit register which, instead of capturing a new value each cycle,
    \textbf{toggles} bit \code{i} only when \code{button_edge(i) = '1'}. Without that gating it would
    toggle every clock cycle instead of once per press, which is the whole point of computing
    \code{button_edge}.
\end{itemize}
This is the whole of \file{led_toggle.vhd}:

\vhdlfile{lectures/L03/led_toggle/led_toggle.vhd}

\textbf{On the board}, as the course demonstrates it: create the Quartus Prime Lite project targeting
the DE0-CV's \code{5CEBA4F23C7N} exactly as in Chapter~\ref{c1:ch} (Appendix~\ref{app:quartus}),
then assign \code{clock} to the \code{50 MHz} oscillator, \code{reset_n} and
\code{button_n(1 downto 0)} to three push buttons, and \code{led(1 downto 0)} to two LEDs. Each
press should toggle exactly one LED, once, no matter how long you hold it.

\textbf{A caveat, deliberately left open:} \code{button_n} is a raw signal straight off a physical
pushbutton pin. In simulation that is a clean, instantaneous transition; on real hardware it is
neither instantaneous (contacts bounce) nor synchronized to \code{clock}. This edge detector is
logically correct but \emph{unsafe} to wire directly to a real button on an FPGA without more work
first, and that work is exactly what Chapter~\ref{c4:ch} covers.

% Section 3.10, from lectures/L03/README.md: the objectives, the questions to test yourself, and
% the next lecture.

\needspace{10\baselineskip}
\section{Review}
\label{c3:sec:review}

\subsection*{What you should be able to do now}
After this chapter, you should be able to:
\begin{itemize}
  \item Say why sequential logic exists, explain the functional difference between a D latch and a
    D flip-flop, and why synchronous designs are built almost entirely from the flip-flop.
  \item Read and write the synchronous process template, and recognize a register as several
    flip-flops sharing a clock, reset and enable.
  \item Say when a signal assignment takes effect: \code{<=} schedules rather than applies, so two
    assignments in one clocked process build a chain, and their order makes no difference.
  \item Build an edge detector from a flip-flop and a gate, and say exactly what signal it produces
    and for how long.
\end{itemize}

\needspace{6\baselineskip}
\subsection*{Questions to test yourself}
You should be able to explain:
\begin{enumerate}
  \item What is the functional difference between a D latch and a D flip-flop, and under what
    condition does each change its output?
  \item Why do synchronous circuits update their state on a single clock edge instead of
    continuously reacting to their inputs?
  \item Given a signal and a flip-flop holding its value from one cycle earlier, how do you combine
    them into a pulse that is high for exactly one clock cycle when the signal rises?
  \item Why must \code{clock} be the only signal in a flip-flop process's sensitivity list, unless
    the design has an asynchronous reset?
  \item A clocked process contains \code{b <= a;} and then \code{c <= b;}. After one clock edge, what
    does \code{c} hold, and why is it not \code{a}? What changes if you swap the two lines, and why?
\end{enumerate}

\needspace{6\baselineskip}
\subsection*{Looking ahead}
Chapter~\ref{c4:ch} is about metastability and synchronization:
\begin{itemize}
  \item Why the raw button feeding this chapter's edge detector can drive a flip-flop into an
    undefined, unstable state, and why that is a different kind of problem from anything here.
  \item The double-flop synchronizer, and what it does and does not do about a bouncing button.
  \item The same LED toggle, rebuilt so it is actually safe on real hardware.
\end{itemize}

% Section 3.11, from lectures/L03/appendix/b_exercises.md.

\section{Exercises}
\label{c3:sec:exercises}

\begin{aside}
\textbf{How to work these.} Each exercise is tagged with the kinds of work it asks for:
\kinds{Paper} by hand, \kinds{CircuitVerse} in the browser, \kinds{VHDL} in a module of your own.

\textbf{How to check your work.} Every exercise that asks you to write a VHDL module ships a
self-checking testbench under \file{lectures/L03/exercises} in the companion repository. Write your
module in its \file{exercises/<module>} directory, using the entity name and the \textbf{port order}
the exercise specifies, then run it with GHDL; \secref{c2:app:c} has the three commands.

No FPGA board is needed for any exercise. Your job is to get the VHDL right, and the testbench is
how you confirm it. Short answers to the parts you work on paper are in
Appendix~\ref{app:answers}; reference solutions to the VHDL are published in the companion
repository, under \file{lectures/L03/solutions}. Try each exercise yourself before looking at
either.
\end{aside}

% ------------------------------------------------------------------------------------------
\exerciseset{D Latch}
\label{c3:set:1}

\exercise{A D latch}{Paper, CircuitVerse}
\label{c3:ex:1}
Build the first circuit in this book that can \emph{remember} something.

Everything you have built so far has been combinational: the outputs follow the inputs, and the
moment the inputs change the outputs change with them. A latch is the smallest departure from that.
It has a data input \code{D} and a control input \code{enable}, and it has two outputs, \code{Q} and
its complement \code{Qn}. While \code{enable} is high the latch is \textbf{transparent}: \code{Q}
simply follows \code{D}, and it looks like a wire. The interesting part is what happens when
\code{enable} goes low. The latch \textbf{locks}, holding whatever \code{D} was at that instant, and
from then on \code{D} can do what it likes without \code{Q} noticing. That held value is one bit of
memory, and it is built out of nothing but the gates you already know, wired so that the output
feeds back into the input.

\code{Qn} is not decoration. The feedback loop is what makes the circuit remember, and the two
outputs are what feed each other: \code{Q} is computed from \code{Qn} and \code{Qn} from \code{Q}.
Build it and you can see the loop; that is the point of drawing it by hand before you ever meet the
VHDL.

Use the equations from \secref{c3:a2}:
\[ Q = (D' \cdot \bool{enable} + \bool{Qn})' \]
\[ \bool{Qn} = (D \cdot \bool{enable} + Q)' \]

\begin{parts}
  \item Realize the corresponding gate network:
    \begin{itemize}
      \item Draw the network by hand.
      \item Recreate and simulate it in CircuitVerse (\url{https://circuitverse.org/simulator}).
    \end{itemize}
  \item Test the transparent and locked states:
    \begin{itemize}
      \item Set \code{enable = 1}:
        \begin{itemize}
          \item Change \code{D}.
          \item Observe what happens to \code{Q} and \code{Qn}.
        \end{itemize}
      \item Set \code{enable = 0}:
        \begin{itemize}
          \item Change \code{D} again.
          \item Observe what happens to \code{Q} and \code{Qn}.
        \end{itemize}
      \item Explain the difference between the two cases.
    \end{itemize}
  \item Cycle \code{D} and \code{enable} through the following combinations twice:
    \begin{itemize}
      \item \code{00}
      \item \code{01}
      \item \code{10}
      \item \code{11}
    \end{itemize}
    Confirm that you observe both cases:
    \begin{itemize}
      \item The latch locks after being transparent with \code{D = 1}.
      \item The latch locks after being transparent with \code{D = 0}.
    \end{itemize}
    Explain how the value of \code{Q} depends on the value of \code{D} at the instant \code{enable}
    changes from \code{1} to \code{0}.
\end{parts}

% ------------------------------------------------------------------------------------------
\exerciseset{D Flip-Flop}
\label{c3:set:2}

\exercise{A D flip-flop from two latches}{Paper, CircuitVerse}
\label{c3:ex:2}
Extend the D latch from Exercise~\ref{c3:ex:1} into a D flip-flop, following \secref{c3:a3}.

The flip-flop must have:
\begin{itemize}
  \item Inputs:
    \begin{itemize}
      \item \code{clock}
      \item \code{reset_n}
      \item \code{D}
      \item \code{enable}
    \end{itemize}
  \item Outputs:
    \begin{itemize}
      \item \code{Q}
      \item \code{Qn}
    \end{itemize}
\end{itemize}
The reset signal is:
\begin{itemize}
  \item Asynchronous:
    \begin{itemize}
      \item It takes effect without waiting for a clock edge.
    \end{itemize}
  \item Active-low:
    \begin{itemize}
      \item Setting \code{reset_n = 0} must force \code{Q} to \code{0} immediately.
    \end{itemize}
\end{itemize}
\note[Tip]
\begin{itemize}
  \item The reset has to do two things in each of the two internal latches, not one:
    \begin{itemize}
      \item Feed an active-high \code{reset} (\code{reset_n} through a NOT gate) into the NOR gate
        that produces \code{Q}. This forces \code{Q} low without waiting for a clock edge.
      \item \code{AND} \code{reset_n} into the AND gate that produces that latch's \code{D} and
        \code{enable} product. This frees \code{Qn} to go high.
      \item Gating only the \code{Q} feedback path is not enough: with \code{enable = 1} and
        \code{D = 1} it leaves \code{Qn} low, and a low \code{Qn} is exactly what holds \code{Q}
        high, so the reset would do nothing.
    \end{itemize}
  \item Drive the enable inputs of the two internal latches using:
    \begin{itemize}
      \item \code{clock}
      \item The inverse of \code{clock}
    \end{itemize}
  \item Follow \secref{c3:a3} for how the external \code{enable} input controls whether the
    flip-flop accepts a new value.
\end{itemize}

\begin{parts}
  \item Realize the corresponding gate network:
    \begin{itemize}
      \item Draw the network by hand.
      \item Recreate and simulate it in CircuitVerse.
      \item Set the clock period to \code{1000 ms}.
        \begin{itemize}
          \item This is slow enough to observe the behaviour by eye.
        \end{itemize}
    \end{itemize}
  \item Test the clocked behaviour:
    \begin{itemize}
      \item Change \code{D} during the high phase of the clock.
      \item Change \code{D} during the low phase of the clock.
      \item Observe \code{Q} and \code{Qn}.
      \item Confirm that the outputs change only on the active clock edge, never between clock
        edges.
    \end{itemize}
    Explain how the two internal latches work together to produce edge-triggered behaviour.
  \item Test the asynchronous reset:
    \begin{itemize}
      \item First set \code{Q = 1}.
      \item Assert the reset by setting \code{reset_n = 0}.
      \item Observe whether \code{Q} resets:
        \begin{itemize}
          \item Immediately.
          \item At the next clock edge.
        \end{itemize}
    \end{itemize}
    Explain the result by referring to the difference between an asynchronous signal and a
    synchronous signal.
\end{parts}

% ------------------------------------------------------------------------------------------
\exerciseset{Registers}
\label{c3:set:3}

\exercise{A four-bit register}{VHDL, Paper}
\label{c3:ex:3}
Write a complete VHDL entity named \code{register4}: four flip-flops side by side, sharing one
clock, one reset, and one enable.

This is the exercise where the flip-flop stops being a curiosity and becomes a building block. A
single D flip-flop stores one bit, which is not much use on its own; what a design actually wants is
somewhere to keep a \emph{number}. Put four flip-flops in a row, wire the same clock to all of them,
and you have a four-bit register. Nothing about the individual flip-flop changes. Only the width
does.

The part worth paying attention to is how little the VHDL has to change to say that. You are
writing the same synchronous process template from \secref{c3:a6} that you already used for one bit,
with \code{std_logic_vector(3 downto 0)} in place of \code{std_logic}. There is no loop, no
repetition, no four copies of anything. That is not a shortcut the language is offering you; it is a
fact about the hardware, and part \textbf{c)} below asks you to say why.

The entity must have:

\begin{booktable}{@{}p{4em}p{4.8em}p{14.6em}L@{}}
  \thead{Port} & \thead{Direction} & \thead{Type} & \thead{Description} \\
  \midrule
  \code{clock} & in & \code{std_logic} & System clock. \\
  \code{reset_n} & in & \code{std_logic} & Asynchronous, active-low. Clears the register to
    \code{"0000"}. \\
  \code{enable} & in & \code{std_logic} & \code{'0'} retains the current value; \code{'1'} captures
    \code{d} on the next rising edge. \\
  \code{d} & in & \code{std_logic_vector(3 downto 0)} & The value to capture. \\
  \code{q} & out & \code{std_logic_vector(3 downto 0)} & The stored value. \\
\end{booktable}

\noindent Write the process out yourself rather than copying the worked D flip-flop example,
\file{d_flip_flop.vhd} (\secref{c3:a8}). Typing it is how the template stops being something you
recognize and starts being something you can produce from memory, and you will be writing it in
every remaining chapter of this book.

\begin{parts}
  \item Write the entity and the synchronous process, adapting the template from a single bit to a
    four-bit vector.
  \item Check it with the testbench, as described at the start of these exercises.
  \item Explain why the same template scales from one flip-flop to four without any structural
    change, while a four-bit \emph{counter} could not simply be four one-bit counters side by side.
    What is different about the two cases?
\end{parts}

\modulefigure{register4}

\selfcheck name your entity \code{register4}, with ports \code{clock}, \code{reset_n},
\code{enable}, \code{d} (\code{std_logic_vector(3 downto 0)}) and \code{q}
(\code{std_logic_vector(3 downto 0)}), declared in that order; its testbench is in
\file{exercises/register4}.

% ------------------------------------------------------------------------------------------
\exerciseset{Edge Detection}
\label{c3:set:4}

\exercise{An edge-detected toggle}{Paper, CircuitVerse}
\label{c3:ex:4}
Construct a circuit that toggles an LED once for every sampled rising edge of a button signal.

Follow:
\begin{itemize}
  \item \secref{c3:a7}
  \item \secref{c3:a9}
\end{itemize}
The circuit must have:
\begin{itemize}
  \item Inputs:
    \begin{itemize}
      \item \code{button}
      \item \code{clock}
      \item \code{reset}
    \end{itemize}
  \item Output:
    \begin{itemize}
      \item \code{led}
    \end{itemize}
\end{itemize}
For this exercise, \code{reset} is active-high, unlike the reset used in the worked example. It is
still \textbf{asynchronous}, like every reset in this book: assert it and \code{led} clears
immediately, without waiting for a clock edge. In the template from \secref{c3:a6} that means the
reset branch stays outside the \code{rising_edge(clock)} test, and only the polarity of the test
changes.

\note[Tip] Use two flip-flops:
\begin{itemize}
  \item One flip-flop stores the previous sampled value of \code{button}.
    \begin{itemize}
      \item Use this value to detect a rising edge.
    \end{itemize}
  \item One flip-flop stores the current state of \code{led}.
    \begin{itemize}
      \item Toggle it only when a rising edge is detected.
    \end{itemize}
\end{itemize}

\begin{parts}
  \item Realize the corresponding gate network:
    \begin{itemize}
      \item Draw the network by hand.
      \item Recreate and simulate it in CircuitVerse.
      \item Set the clock period to \code{1000 ms}.
    \end{itemize}
  \item Test the edge detector:
    \begin{itemize}
      \item Hold \code{button} high for several consecutive clock cycles.
      \item Confirm that \code{led} toggles exactly once.
      \item Confirm that it does not toggle once per clock cycle.
    \end{itemize}
    Explain which stored signal prevents the circuit from repeatedly detecting the same high level
    as a new rising edge.
\end{parts}

\exercise{The toggle in VHDL}{VHDL}
\label{c3:ex:5}
Implement the circuit from Exercise~\ref{c3:ex:4} in VHDL.

\begin{parts}
  \item Create an entity named \code{led_toggle_single} with:

\begin{booktable}{@{}p{4em}p{4.8em}p{5.2em}L@{}}
  \thead{Port} & \thead{Direction} & \thead{Type} & \thead{Description} \\
  \midrule
  \code{clock} & in & \code{std_logic} & System clock. \\
  \code{reset} & in & \code{std_logic} & Asynchronous, and \textbf{active-high}, unlike every other
    module in this book. Assert it and \code{led} clears without waiting for a clock edge, so the
    reset branch belongs outside the \code{rising_edge(clock)} test. \\
  \code{button} & in & \code{std_logic} & Toggles \code{led} once per press. \\
  \code{led} & out & \code{std_logic} & The LED state. \\
\end{booktable}

    Note the single-bit types: the worked example uses \code{std_logic_vector} ports for two buttons
    and two LEDs, and this one drives a single pair.

    Do not copy \file{led_toggle.vhd} (\secref{c3:a9}) directly. Instead:
    \begin{itemize}
      \item Re-derive the two clocked processes yourself:
        \begin{itemize}
          \item One process stores the previous value of \code{button}.
          \item One process stores and toggles the current value of \code{led}.
        \end{itemize}
      \item Derive the edge-detection equation for:
        \begin{itemize}
          \item A rising edge.
          \item An active-high button.
        \end{itemize}
      \item Use the rising-edge equation from \secref{c3:a7}:
        \[ \bool{edge} = \bool{current} \cdot \bool{previous}' \]
    \end{itemize}
    Do not use the falling-edge equation from the worked example, which uses an active-low button.
  \item Verify the design with its testbench (see the aside at the start of these exercises).

    The testbench drives the same sequence you would exercise by hand on a board:
    \begin{itemize}
      \item It asserts \code{reset} and checks \code{led} is cleared.
      \item It raises \code{button} and checks \code{led} toggles exactly once.
      \item It \textbf{holds \code{button} high} for several more clock cycles and checks, after
        every one of them, that \code{led} does \emph{not} toggle again.
      \item It releases and presses again, and checks the second press toggles \code{led} once more:
        back off, since the first press had turned it on.
      \item It releases and presses a third time, so the LED is lit when the reset arrives.
      \item It asserts \code{reset} with \code{led} lit and \textbf{no clock edge following}, and
        checks \code{led} clears anyway. This is the one that fails if you put the reset inside the
        \code{rising_edge(clock)} branch.
    \end{itemize}
    That third check is the one that fails if your edge detector is wrong:
    \begin{itemize}
      \item An \code{led} that toggles on every clock cycle while the button is held means you drove
        the toggle from \code{button}'s \textbf{level} rather than from its edge.
      \item Re-read the equation above: \code{edge} is high only when \code{button} is high
        \emph{now} and its stored previous value is low.
    \end{itemize}
\end{parts}

\modulefigure{led_toggle_single}

\selfcheck name your entity \code{led_toggle_single}, with ports \code{clock}, \code{reset},
\code{button} (in) and \code{led} (out), declared in that order; its testbench is in
\file{exercises/led_toggle_single}. Note \code{reset} is \textbf{active-high} here, unlike the
worked example.

